%% wave12_a10_wall_crossing_soibelman.tex
%% Wave 12 / DNA strand A10: adversarial assault and rigorous healing of the
%% wall-crossing / Maurer--Cartan gauge-equivalence claims, the conifold MC
%% chamber bridge, the grand atlas of CY3 families, and the CY4 prediction.
%% Soibelman voice: DT-categorical rigour, CoHA precision, KS wall-crossing
%% vigilance. Reader-facing target sections:
%%   working_notes.tex 410--428  (Wall-crossing and MC gauge equivalence)
%%   working_notes.tex 1429--1453 (Conifold: wall-crossing as MC gauge equivalence)
%%   working_notes.tex 1541--1591 (Grand atlas of CY3 families)
%%   working_notes.tex 1792--1817 (Updated open questions)
%%   working_notes.tex 2527--2588 (CY_4 prediction)
%% This file is infrastructure (\texttt{notes/}), not manuscript prose; it
%% states the mathematics in full, records ghost-theorem/error/correction
%% for each attack, and specifies the exact edits required for Vol~III rectification.

\documentclass[11pt]{article}
\usepackage[margin=1.15in]{geometry}
\usepackage{amsmath, amssymb, amsthm, mathtools}
\usepackage{enumitem}
\usepackage{hyperref}
\usepackage{xcolor}
\usepackage{booktabs}
\usepackage{longtable}

\newtheorem{theorem}{Theorem}
\newtheorem{proposition}[theorem]{Proposition}
\newtheorem{lemma}[theorem]{Lemma}
\newtheorem{corollary}[theorem]{Corollary}
\newtheorem{conjecture}[theorem]{Conjecture}
\theoremstyle{definition}
\newtheorem{definition}[theorem]{Definition}
\newtheorem{remark}[theorem]{Remark}
\newtheorem{observation}[theorem]{Observation}
\newtheorem{attack}{Attack}
\newtheorem{heal}{Heal}

\providecommand{\C}{\mathbb{C}}
\providecommand{\Z}{\mathbb{Z}}
\providecommand{\R}{\mathbb{R}}
\providecommand{\Q}{\mathbb{Q}}
\providecommand{\N}{\mathbb{N}}
\providecommand{\bP}{\mathbb{P}}
\providecommand{\bS}{\mathbb{S}}
\providecommand{\bR}{\mathbb{R}}
\providecommand{\cO}{\mathcal{O}}
\providecommand{\cM}{\mathcal{M}}
\providecommand{\cF}{\mathcal{F}}
\providecommand{\cA}{\mathcal{A}}
\providecommand{\cH}{\mathcal{H}}
\providecommand{\CoHA}{\mathrm{CoHA}}
\providecommand{\Stab}{\mathrm{Stab}}
\providecommand{\Hall}{\mathrm{Hall}}
\providecommand{\Eone}{E_1}
\providecommand{\Etwo}{E_2}
\providecommand{\Einf}{E_\infty}
\providecommand{\ad}{\mathrm{ad}}
\providecommand{\BCH}{\mathrm{BCH}}
\providecommand{\kch}{\kappa_{\mathrm{ch}}}
\providecommand{\kcat}{\kappa_{\mathrm{cat}}}
\providecommand{\kBKM}{\kappa_{\mathrm{BKM}}}
\providecommand{\kfib}{\kappa_{\mathrm{fiber}}}

\title{Wall-crossing as Maurer--Cartan gauge equivalence \\ {\normalfont (Soibelman-voice adversarial rectification, strand A10)}}
\author{Raeez Lorgat}
\date{\today}

\begin{document}
\maketitle

\begin{abstract}
Five adversarial assaults on the \emph{wall-crossing $=$ MC gauge
equivalence} programme in working\_notes.tex, each followed by a
first-principles healing. The conifold test case is reconstructed from
the Kontsevich--Soibelman motivic Hall algebra downward: the
semi-classical dgLa $L_\Gamma$, the two chamber MC elements
$\Theta_\mathrm{I}, \Theta_\mathrm{II}$, and the explicit chamber-
interpolating gauge $\lambda \in L_\Gamma^0$ written as a path-ordered
integral over a Kähler-moduli wall crossing. The grand atlas is
rectified (class assignments, $\kch$ entries, conjecture-vs-theorem
status, the conifold's $d=3$ direct-McKay invariant $\kch = 1$ against
AP-CY6). The CY$_4$ heuristic is sharpened. AP-CY-WC1 through AP-CY-WC5
are appended.
\end{abstract}

%======================================================================
\section{Gate 0: context and posture}
%======================================================================

\paragraph{Primary sources.}
Kontsevich--Soibelman, \emph{Stability structures, motivic DT-invariants and
cluster transformations} (\texttt{arXiv:0811.2435}); Kontsevich--Soibelman,
\emph{Cohomological Hall algebra, exponential Hodge structures and motivic
DT-invariants} (\texttt{arXiv:1006.2706}); Szendr\H{o}i, \emph{Non-commutative
Donaldson--Thomas theory of the conifold} (\emph{Geom.\ Topol.}\ 12, 2008);
Morrison--Mozgovoy--Nagao--Szendr\H{o}i, \emph{Motivic DT invariants of the
conifold and an ADE generalisation} (\texttt{arXiv:1010.2354}); Davison
\texttt{arXiv:1209.4620}; Nagao \texttt{arXiv:0907.3784}; Reineke
\emph{The use of geometric and quantum group techniques for wild quivers}
(\texttt{math/0404054}); Gross--Siebert \texttt{arXiv:0911.2780}; Bridgeland
\emph{Stability conditions on triangulated categories} (\texttt{math/0212237}).

\paragraph{Cache posture (CLAUDE.md authority).}
\begin{itemize}[leftmargin=1.1em]
\item CY-C is \emph{conjectural} (AP-CY6, top-15 \#10). No theorem stated in
  this note relies on CY-C at $d = 3$.
\item Unsubscripted $\kappa_\bullet$ is forbidden (AP113, HZ-7); every
  $\kappa_\bullet$ symbol below carries a subscript in
  $\{\mathrm{ch}, \mathrm{cat}, \mathrm{BKM}, \mathrm{fiber}\}$.
\item Conifold is \emph{not} a local surface at $d = 3$; direct McKay gives
  $\kch(\text{conifold}) = 1$. The atlas entry must cite direct McKay or
  DT genus-$1$, not a dimensional reduction.
\item $\CoHA(\C^3) = Y^+$ (positive half), \emph{not} $\mathcal{W}_{1+\infty}$.
\item AP-CY57 (construction/narration needs explicit arrow): every gauge
  equivalence below names the element $\lambda \in L^0$, the one-parameter
  family $\{m_t\}_{t \in [0,1]}$, and the explicit vector field generating
  the gauge action.
\end{itemize}

%======================================================================
\section{Gauge data: the dgLa behind wall-crossing}
\label{sec:dgla}
%======================================================================

Fix a $\Z$-graded abelian category $\cA$ with a stability structure
$\sigma \in \Stab(\cA)$ (in the sense of Bridgeland). Let
$\Gamma = K_0(\cA)$ and fix the skew-symmetric Euler pairing
$\langle \cdot, \cdot \rangle\colon \Gamma \times \Gamma \to \Z$ induced by
$\chi_\cA - \chi_\cA^{\mathrm{op}}$.

\begin{definition}[Lattice algebra and motivic quantum torus]
\label{def:lattice-algebra}
The \emph{lattice Lie algebra} $\mathfrak{g}_\Gamma$ is the
$\Q$-vector space spanned by formal symbols
$\{ e_\gamma : \gamma \in \Gamma \setminus \{0\} \}$ with bracket
\[
[e_{\gamma_1}, e_{\gamma_2}] \;=\; (-1)^{\langle \gamma_1, \gamma_2 \rangle}
\langle \gamma_1, \gamma_2 \rangle \, e_{\gamma_1 + \gamma_2}.
\]
The \emph{quantum torus} $\mathfrak{g}_\Gamma^q$ replaces the
$(-1)^{\langle \cdot, \cdot \rangle}$ sign by $q^{\frac{1}{2}
\langle \cdot, \cdot \rangle}$ and the bracket by
$[\hat e_{\gamma_1}, \hat e_{\gamma_2}] = (q^{\langle \gamma_1, \gamma_2
\rangle/2} - q^{-\langle \gamma_1, \gamma_2 \rangle/2}) \hat e_{\gamma_1 +
\gamma_2}$. In the semi-classical limit $q \to 1$ the second formula
recovers the first up to the KS $(-1)$-twist (Kontsevich--Soibelman
\texttt{0811.2435}, \S2.5).
\end{definition}

\begin{definition}[Central-charge filtration and the pro-nilpotent dgLa]
\label{def:dgla}
Let $Z\colon \Gamma \to \C$ be a central charge. For each open strict convex
cone $V \subset \bR^2$ set
\[
\mathfrak{g}_\Gamma[V] \;=\; \widehat{\bigoplus}_{\gamma: Z(\gamma) \in V}
\Q \cdot e_\gamma,
\]
the completion taken along $|Z(\gamma)| \to \infty$. Equip
$\mathfrak{g}_\Gamma[V]$ with zero differential. This is the
\emph{semi-classical KS dgLa}. It is concentrated in degree zero; its MC
equation is trivially $[m, m]/2 = 0$, so \emph{every} element is an MC
element. The nontrivial data are the chamber-specific MC elements
\[
\Theta_V \;=\; \sum_{\gamma:\, Z(\gamma) \in V} \Omega(\gamma) \, e_\gamma,
\qquad \Omega(\gamma) \in \Z,
\]
where $\Omega(\gamma) = \Omega(\gamma; \sigma)$ are the numerical
DT-invariants of $\sigma$-semistable objects of class $\gamma$.
\end{definition}

\begin{remark}[Lie-theoretic vs Hall-algebra lift]
\label{rem:lift}
$\mathfrak{g}_\Gamma[V]$ is the \emph{Lie-theoretic} wall-crossing
target. The Hall-algebra lift replaces $e_\gamma$ by the integration-map
class of the moduli stack $\cM^\sigma_\gamma$ in the motivic Hall
algebra $H(\cA)$, and $\Omega(\gamma)$ is recovered by the Reineke
stability-sum formula (\texttt{math/0404054}). The wall-crossing
identity is an identity of group-like elements in $\widehat{H(\cA)}$
which, after applying the integration map, descends to an identity in
$\exp(\mathfrak{g}_\Gamma[V])$.
\end{remark}

\begin{theorem}[KS wall-crossing as MC gauge equivalence; semi-classical]
\label{thm:ks-wc-mc}
Let $\sigma_0, \sigma_1 \in \Stab(\cA)$ lie in two adjacent chambers
separated by a single active wall $W$, and let $V_0, V_1 \subset \bR^2$
be half-plane cones just past $W$ for $\sigma_0, \sigma_1$ respectively.
Write $\Theta_0 = \Theta_{V_0}, \Theta_1 = \Theta_{V_1}$. Then there
exists $\lambda \in \mathfrak{g}_\Gamma^{W} \subset \mathfrak{g}_\Gamma$
supported on the single-wall cone $W \cdot \R_{>0}$ with
\[
e^{\ad_\lambda}(\Theta_0) \;=\; \Theta_1.
\]
The element $\lambda$ is determined by the KS factorisation of the single
wall product
\[
\prod^{\curvearrowleft}_{\gamma \in W} T_\gamma^{\Omega(\gamma)}
\;=\; \prod^{\curvearrowright}_{\gamma \in W} T_\gamma^{\Omega'(\gamma)}
\]
through the log map, where $T_\gamma = \exp(e_\gamma)$ and the products
are taken in the profinite group $\exp(\mathfrak{g}_\Gamma[W \cdot \R_{>0}])$.
The DT invariants $\{\Omega, \Omega'\}$ for the two chambers are
determined by the wall crossing; their difference is the
wall-crossing formula (\texttt{0811.2435}, Thm.~5).
\end{theorem}

\begin{proof}[Proof sketch]
The wall $W$ partitions the central-charge plane into a phase $\phi_W$;
the charges $\gamma$ with $Z_{\sigma_0}(\gamma), Z_{\sigma_1}(\gamma)$
straddling $W$ are exactly those with $Z(\gamma) \in \R_{>0} W$. The
product $T_W^{\sigma_0} = \prod_\gamma T_\gamma^{\Omega_{\sigma_0}(\gamma)}$
defines a single automorphism $g_W$ of the pro-unipotent group
$\exp(\mathfrak{g}_\Gamma)$ by taking its log. Set
$\lambda = \log(g_W)$; its support is contained in $W \cdot \R_{>0}$
because $\log$ is a convergent BCH series in the pro-nilpotent regime
(the cone condition ensures BCH convergence). The KS wall-crossing
identity $T_W^{\sigma_0} = T_W^{\sigma_1}$ means $g_W$ reads identically
from either chamber. Then
$e^{\ad_\lambda}(\Theta_0) = \mathrm{Ad}_{g_W}(\Theta_0)$, which by a
direct computation using the action of $T_\gamma$ by conjugation
(shifting charges by $\gamma$) equals $\Theta_1$ precisely when
$\Omega_{\sigma_0}, \Omega_{\sigma_1}$ are related by the KS formula.
\end{proof}

\paragraph{What this gives and what it does not.}
Theorem~\ref{thm:ks-wc-mc} is the \emph{semi-classical} identification
used in working\_notes \S\ref{wn:sec:wall-crossing-mc} (conifold,
pentagon). At the motivic (quantum) level, MC gauge equivalence becomes
a conjugation in $\exp(\mathfrak{g}_\Gamma^q)$ and the quantum dilogarithm
$E(x) = \prod_{k\ge 0}(1 + q^{k+1/2} x)^{-1}$ replaces $T_\gamma$
(\texttt{0811.2435}, \S6). The pentagon $E(x)E(y) = E(y)E(xy)E(x)$ for
$yx = q xy$ is the quantum conifold identity.

%======================================================================
\section{Attacks and heals}
%======================================================================

\subsection*{Attack 1: Is ``wall-crossing as MC gauge equivalence'' an
abuse of notation for chamber change in $\Stab(\cA)$?}

\begin{attack}[\textbf{Ghost theorem} / \textbf{precise error} / \textbf{correct relation}]
\label{atk:1}
The line 414--416 remark asserts that the KS wall-crossing formula ``is
verified as a gauge equivalence of Maurer--Cartan elements''. A reader
testing this against the standard deformation-theory meaning of MC
gauge equivalence---a dgLa $L$ with MC set $\mathrm{MC}(L) = \{ m \in
L^1 : dm + \tfrac{1}{2}[m,m] = 0\}$ modulo $\exp(L^0)$---finds
the dgLa unnamed, $L^1$ unidentified, $\lambda \in L^0$ unspecified,
and $\Theta$ appearing without a listed underlying complex.

\textbf{Ghost theorem}: the identification is real and reduces to
Theorem~\ref{thm:ks-wc-mc}.

\textbf{Precise error}: the remark elides the three pieces of data
(which dgLa, which MC element, which gauge). Without naming them, the
statement is Pattern AP-CY57 (construction/narration without explicit
arrow). The pentagon identity $\BCH(L_{10}, L_{01}) = \BCH(L_{01}, L_{11},
L_{10})$ is verified through height 2, but height 3+ is exactly where
the Gross--Siebert scattering-diagram obstruction (scattering diagram
$=$ $\Eone$ MC gauge at motivic-Hall-algebra level but not naive-Lie
level; working\_notes \S2547) appears. So the raw BCH pentagon is a
low-height fact; the genuine identification is with the Hall-algebra
integration and requires the motivic lift.

\textbf{Correct relation}: the dgLa is $\mathfrak{g}_\Gamma[V]$ of
Def.~\ref{def:dgla} (concentrated in degree 0; trivial differential);
MC elements are $\Theta_V = \sum_\gamma \Omega_V(\gamma) e_\gamma$;
$\lambda$ is the log of the single-wall motivic product. The pentagon
is true \emph{to all orders} in the quantum-torus dgLa, but only to
height 2 in the Lie-theoretic semi-classical dgLa with
$(-1)^{\langle \cdot, \cdot \rangle}$ signs; the height-$\ge 3$
correction is supplied by the Reineke integration map on the
motivic side, not by a naive BCH expansion.
\end{attack}

\begin{heal}[dgLa, $\Theta$, $\lambda$ explicitly named]
\label{heal:1}
Replace the working\_notes remark (lines 414--416) with the following
direct statement. The replacement parses as a theorem with explicit
arrow, discharging AP-CY57.

\smallskip
\textit{Replacement text.}
For the resolved conifold $\mathbb{Y}$, Theorem~\ref{thm:ks-wc-mc}
specialises to the following. Take $\Gamma = \Z\gamma_1 \oplus \Z\gamma_2$
with $\langle \gamma_1, \gamma_2 \rangle = 1$; fix the wall $W$ at
phase $\phi_W$ where $Z(\gamma_1) \parallel Z(\gamma_2)$. Write the
chamber-$\mathrm{I}$ MC element
\[
\Theta_\mathrm{I} \;=\; e_{\gamma_1} + e_{\gamma_2},
\]
and the chamber-$\mathrm{II}$ MC element
\[
\Theta_\mathrm{II} \;=\; e_{\gamma_1} + e_{\gamma_1 + \gamma_2} + e_{\gamma_2},
\]
recording the $\Omega(\gamma_1) = \Omega(\gamma_2) = \Omega(\gamma_1 + \gamma_2) = 1$
DT invariants past the wall.
The gauge $\lambda \in \mathfrak{g}_\Gamma$ is
\[
\lambda \;=\; -\tfrac{1}{2} [e_{\gamma_1}, e_{\gamma_2}]
\;=\; -\tfrac{1}{2} \, (-1)^{\langle \gamma_1, \gamma_2 \rangle}
\langle \gamma_1, \gamma_2 \rangle \, e_{\gamma_1 + \gamma_2}
\;=\; \tfrac{1}{2} e_{\gamma_1 + \gamma_2},
\]
(sign $(-1)^{1} = -1$; $\langle \gamma_1, \gamma_2 \rangle = 1$). Then
$e^{\ad_\lambda}(\Theta_\mathrm{I}) = \Theta_\mathrm{II}$ through
height 2 (BCH truncation at two nested brackets) in the Lie-theoretic
dgLa, and to all orders in the quantum-torus dgLa via
$E(e_{\gamma_1}) E(e_{\gamma_2}) = E(e_{\gamma_2}) E(e_{\gamma_1 +
\gamma_2}) E(e_{\gamma_1})$. \hfill $\square$

\smallskip
\textit{Audit status.} The Reineke $\Omega = 1$ convention is used
throughout; the fermionic convention $\Omega = -1$ does not satisfy
the pentagon identity with the same wall ordering (as noted in
working\_notes line 416).
\end{heal}

\subsection*{Attack 2: Is the conifold flop a gauge equivalence, and in
which dgLa?}

\begin{attack}
\label{atk:2}
Working\_notes line 1438--1442 writes
$\exp(\ad_\alpha)(\Theta_\mathrm{I}) = \Theta_\mathrm{II}$ ``where
$\alpha$ lies in the lattice Lie algebra''. Two distinct notions of
``chamber'' coexist in the conifold literature:
\begin{enumerate}[label=(\alph*), leftmargin=1.1em]
\item \emph{Kähler chambers}: the resolved conifold
  $\cO(-1) \oplus \cO(-1) \to \bP^1$ and its flop. These chambers
  are connected by a birational transformation preserving $\kch$
  (AP-CY10: flop $\ne$ Koszul dual).
\item \emph{Stability chambers in $\Stab(D^b(\bY))$}: the
  Szendr\H{o}i non-commutative DT chamber (where the Jacobi algebra
  $J(Q_\mathrm{con}, W)$ is abelian-tilted), the NCDT chamber of
  Mozgovoy--Reineke, and the large-volume DT chamber.
\end{enumerate}
\textbf{Ghost theorem}: both notions are instances of KS wall-crossing,
but on \emph{different} walls in the larger moduli.

\textbf{Precise error}: the working\_notes conflates (a) and (b) by
writing ``the two DT chambers (large-volume and flopped)'' (line 1432).
The flopped resolution is a different CY3 with its own $\Stab$; the
flop is a birational wall in Kähler moduli, not a stability wall in a
fixed $\Stab(D^b(\bY))$. The relevant wall for $\CoHA(\bY)$ is the
Szendr\H{o}i NCDT $\leftrightarrow$ large-volume DT wall on a fixed
triangulated category.

\textbf{Correct relation}: there are two towers of walls. The
\emph{birational} (flop) wall is detected by Bridgeland's derived
equivalence $\Phi^{\mathrm{flop}}\colon D^b(\bY) \to D^b(\bY^+)$, which
induces an isomorphism on the lattice $\Gamma$ and on
$\mathfrak{g}_\Gamma$. The \emph{stability} walls inside a fixed
$\Stab(D^b(\bY))$ are the ones governed by Theorem~\ref{thm:ks-wc-mc}.
The pentagon identity of working\_notes \S\ref{wn:sec:wall-crossing-mc}
is a stability wall, not a flop wall.
\end{attack}

\begin{heal}[Two-wall stratification on the conifold]
\label{heal:2}
State both walls, cleanly separated, and rename the working\_notes
subsection's ``(large-volume and flopped)'' (line 1432) to
``(Szendr\H{o}i NCDT and large-volume DT)''.

\smallskip
\textit{Stability wall (governing Theorem~\ref{thm:conifold-wc}).}
In $\Stab(D^b(\bY))$ there is a distinguished wall $W_\mathrm{DT}$
where the central charge of $\cO_\bY[1]$ and the sum
$\cO_{\mathbb{P}^1} + \cO_{\mathbb{P}^1}(-1)[1]$ align. On one side
(Szendr\H{o}i's NCDT chamber), the simple objects of the abelian heart
are the two $Q_\mathrm{con}$-simples $S_\circ, S_\bullet$ of dimension
vectors $(1, 0), (0, 1)$; on the other side (large-volume DT), they
are $\cO_{\mathrm{pt}}$ and $\cO_{\bP^1}(-1)[1]$. The MC elements are
\[
\Theta_\mathrm{NCDT} \;=\; e_{\gamma_\circ} + e_{\gamma_\bullet},
\qquad
\Theta_\mathrm{LV} \;=\; e_{\gamma_\circ} + e_{\gamma_\circ + \gamma_\bullet}
+ e_{\gamma_\bullet}.
\]
Heal~\ref{heal:1} applies verbatim with $\gamma_1 = \gamma_\circ$,
$\gamma_2 = \gamma_\bullet$.

\smallskip
\textit{Birational wall (flop).}
The flop $\bY \dashrightarrow \bY^+$ induces the derived equivalence
$\Phi^{\mathrm{flop}}$ of Bridgeland--King--Reid, and the CoHA and
$\kch$ are preserved (AP-CY10). On the lattice:
$\Phi^{\mathrm{flop}}(\gamma_\circ) = -\gamma_\bullet$,
$\Phi^{\mathrm{flop}}(\gamma_\bullet) = \gamma_\circ + \gamma_\bullet$
(Atiyah flop formula). This is \emph{not} captured by
Theorem~\ref{thm:ks-wc-mc}; it is a separate identification and must
be stated as such.

\smallskip
\textit{Audit status.} The working\_notes wording
``flopped resolution'' (line 1432) is misleading; the correct wording
is ``Szendr\H{o}i NCDT and large-volume DT chambers on a fixed
triangulated category $D^b(\bY)$, related by the stability wall
$W_\mathrm{DT}$''. The flop relation $\bY \leftrightarrow \bY^+$ is a
\emph{separate} equivalence (Atiyah flop).
\end{heal}

\subsection*{Attack 3: Is $\kch(\text{conifold}) = 1$ from ``DT genus 1'' the
same as from McKay?}

\begin{attack}
\label{atk:3}
Working\_notes line 1437 asserts $\kch(\text{conifold}) = 1$ with class
$\mathbf{G}$ and $r_{\max} = 2$. The atlas (line 1560) attributes this
to ``DT genus-1''.

\textbf{Ghost theorem}: $\kch(\text{conifold}) = 1$ is a genuine
Vol~III constant, cached in CLAUDE.md (``Conifold is NOT a local
surface at $d=3$; $\kch = 1$ via direct McKay'').

\textbf{Precise error}: ``DT genus 1'' is ambiguous---does this refer
to (i) the genus-1 Gopakumar--Vafa free energy
$F_1^{\mathrm{con}} = -\tfrac{1}{12}\log(1 - Q)$, (ii) a genus-1 bar-
complex computation, or (iii) genus-1 modular behaviour of the DT
partition function? The cache demands \emph{direct McKay}: the
conifold is the $\C^\times$-gerbe over the local affine plane with the
Szendr\H{o}i quiver as McKay data, and $\kch = 1$ reads off the
quiver's vertex count minus one relation. The conifold is \emph{not} a
local surface (it is genuinely $d = 3$, total space is $6$-real-
dimensional and $h^{1,1} = 1$ reflects the resolving $\bP^1$, not a
surface structure). Calling this ``DT genus 1'' invites conflation with
$F_1^{\mathrm{con}}$, an \emph{unrelated} numerical observable.

\textbf{Correct relation}: $\kch(\text{conifold}) = 1$ via the direct
McKay computation: $\CoHA(\bY) \cong U(\mathfrak{g}_\mathrm{BPS}(\bY))$
where $\mathfrak{g}_\mathrm{BPS}(\bY)$ is a $\Z^2$-graded Lie algebra
with a single rank-1 generator at each primitive charge
$\gamma \in \{\gamma_\circ, \gamma_\bullet, \gamma_\circ + \gamma_\bullet\}$;
the shadow-tower computation collapses to a rank-1 Gaussian at all
primitive charges, giving $\kch = 1$. The ``genus 1'' of the free
energy $F_1^{\mathrm{con}}$ is a topological-string observable and
distinct from $\kch$.
\end{attack}

\begin{heal}[Atlas source column corrected]
\label{heal:3}
Change the atlas row for the conifold (working\_notes line 1560) from
\begin{quote}
\texttt{Conifold \& --- \& 1 \& 0 \& 1 \& G \& DT genus-1}
\end{quote}
to
\begin{quote}
\texttt{Conifold \& --- \& 1 \& 0 \& 1 \& G \& direct McKay ($\CoHA(\bY)$ rank-1 per charge)}
\end{quote}
The source column becomes consistent with the CLAUDE.md cache key
``Conifold is NOT a local surface at $d = 3$; $\kch = 1$ via direct
McKay''. Cross-reference: \texttt{notes/CoHA\_to\_W\_infty\_treatise.tex}
Example 2 (\S4), which gives the direct McKay chain.
\end{heal}

\subsection*{Attack 4: Is the grand atlas row for $K3 \times E$ missing the
four-fold $\kappa_\bullet$ spectrum?}

\begin{attack}
\label{atk:4}
Working\_notes line 1563 writes for $K3 \times E$:
``$\kch = 3$ ($\kBKM = 5$)''. This is \emph{half} the canonical
spectrum. The CLAUDE.md constant says
``K3 $\times$ E spectrum $\{2, 3, 5, 24\}$ (four distinct
constructions)'': Mukai lattice ($\kcat = 2$ at K3 fibre only;
Künneth-multiplicative $\kcat(K3 \times E) = 0$ on total space);
Igusa $\Phi_{10}$ via Gritsenko $\Delta_5$ gives
$\kBKM = \mathrm{wt}(\Delta_5) = 5$; BKM Borcherds weight gives $5$;
K3 fibre-rank 24 gives $\kfib = 24$; and $\kch = 3$ is the
Euler-supertrace on the total space.

\textbf{Ghost theorem}: the atlas row is \emph{telling one face} of the
four-face spectrum, not the full spectrum.

\textbf{Precise error}: omitting $\kfib$ and $\kcat$ from the $K3 \times
E$ row invites readers to infer a $\kBKM = \kch + \kfib$ identity or
confuse $\kch(K3 \times E) = 3$ with $\kcat(K3 \times E)$. The
CLAUDE.md cache explicitly says $\kBKM = \kch + \chi(\cO_\mathrm{fiber})$
is \emph{false} at every $N \in \{1, 2, 3, 4, 6\}$ (at $N=1$: left $= 5$,
right $= 0 + 0 = 0$). The correct identity is $\kBKM(\Phi_N) = c_N(0)/2$.

\textbf{Correct relation}: the $K3 \times E$ row is the Vol~III
crystallisation point. Its four entries $\{\kch, \kcat, \kBKM, \kfib\}
= \{3, 0, 5, 24\}$ come from four different constructions. The
atlas should record all four.
\end{attack}

\begin{heal}[Atlas $K3 \times E$ row expanded]
\label{heal:4}
Replace the working\_notes $K3 \times E$ row with
\begin{center}
\small
\begin{longtable}{lccccl}
\toprule
CY3 & $\chi$ & $h^{1,1}$ & $h^{2,1}$ & $\{\kch, \kcat, \kBKM, \kfib\}$ & Source \\
\midrule
$K3 \times E$ & $0$ & $21$ & $21$ & $\{3, 0, 5, 24\}$ &
\begin{tabular}{@{}l@{}}$\kch$: Hodge supertrace\\
$\kcat$: Künneth, $\chi(\cO_{K3}) \chi(\cO_E) = 2\cdot 0$\\
$\kBKM = c_1(0)/2$: Gritsenko $\Delta_5$, weight $5$\\
$\kfib$: K3 fibre-rank $= \chi(K3) = 24$\end{tabular} \\
\bottomrule
\end{longtable}
\end{center}
Replace the single-line observation ``(iv) Only 9 of the 14+ catalogued
families have integer $\kch$; a genuine BKM algebra requires $\kch \in
\Z_{>0}$'' (line 1583) by:
\begin{quote}
\emph{(iv)} The K3-BKM family $\{K3 \times E, \text{Enriques} \times E,
\text{Borcea--Voisin (K3 type), }\ldots\}$ has $\kBKM = c_N(0)/2$
universal (Borcherds 1995, Gritsenko 1999); the table records $\kBKM$
only where the relevant Siegel-modular denominator form is known. A
\emph{distinct} line of families has $\kch = \chi_\mathrm{top}/24 \in \Q$
(quintic, $\bP^5[3,3]$, $\bP^5[2,4]$, $\ldots$); these do not carry a
$\kBKM$ without an auxiliary automorphic construction.
\end{quote}
The six-routes-to-$G(K3\times E)$ discussion is a companion to the
$\{\kch, \kcat, \kBKM, \kfib\}$ spectrum but concerns a different
claim (six different constructions of one object); it should not be
silently merged with the four $\kappa_\bullet$ values.
\end{heal}

\subsection*{Attack 5: Is the open question ``wall-crossing as MC gauge
equivalence: verified for the conifold'' actually \emph{resolved}?}

\begin{attack}
\label{atk:5}
Working\_notes line 1802 upgrades this question to \emph{verified for
the conifold} citing Theorem~\ref{thm:conifold-wc}. The remaining
open part (general CY3) is not stated.

\textbf{Ghost theorem}: Theorem~\ref{thm:ks-wc-mc} \emph{is} the
general statement and holds for any CY3 with a reasonable stability
structure, subject to caveats.

\textbf{Precise error}: the open question is not closed; its scope has
narrowed. Caveats include: (i) the semi-classical identification
requires the Lie-theoretic dgLa $\mathfrak{g}_\Gamma$, but
identifications at the $\Etwo$-chiral level (see the $E_n$ hierarchy
Part of Vol~III) require
lifting to the motivic Hall algebra $H(\cA)$; (ii) for non-toric CY3
(quintic, Gepner chart), the stability manifold $\Stab(D^b(X))$ is
poorly understood, and the KS integrality conjecture for $\Omega$ is
open (\texttt{1006.2706}, \S7); (iii) at $d = 3$ there is no $A_X$
(AP-CY6), so the identification ``wall-crossing = MC gauge of
$A_X$'' is ill-posed; the correct formulation is ``wall-crossing $=$
MC gauge of $\CoHA(X)$ at the Hall-algebra level''.

\textbf{Correct relation}: the open question (iv) of
working\_notes line 1802 should be revised to read: ``Wall-crossing as
MC gauge equivalence: \emph{established as a theorem}
(Theorem~\ref{thm:ks-wc-mc}) at the semi-classical KS dgLa for any
CY3 admitting a stability structure; verified explicitly on the
conifold (Theorem~\ref{thm:conifold-wc}) and on $\C^3$ (MNOP
DT/PT wall). The $\Etwo$-chiral refinement and the motivic-Hall-
algebra lift remain open at $d = 3$ pending CY-A$_3$ morphism-level
results (AP-CY6, CY-C conjectural).''
\end{attack}

\begin{heal}[Open-question item 4 rewritten]
\label{heal:5}
Replace working\_notes line 1802 by
\begin{quote}
\textbf{Wall-crossing as MC gauge equivalence}: \emph{established} at
the semi-classical KS dgLa for any CY3 admitting a stability structure
(Theorem~\ref{thm:ks-wc-mc}); verified explicitly on the conifold
(Theorem~\ref{thm:conifold-wc}) and on $\C^3$. The $\Etwo$-chiral
refinement, and the motivic Hall-algebra lift relating MC gauge to
the Reineke integration map, remain open at $d = 3$ pending
CY-A$_3$ morphism-level results (AP-CY6; CY-C conjectural).
\end{quote}
Companion open-question edits:
\begin{itemize}[leftmargin=1.1em]
\item Item 5 ($\kch = h^{1,1}/4$ universal?): \emph{No} is correct;
  append ``and in general $\kch$ is the $\Eone$-chiral Euler
  supertrace $\sum_q (-1)^q h^{0,q}(X)$ on compact CY$_d$, not a
  function of $h^{1,1}$ alone.''
\item Item 13 (Shadow--BPS/DT correspondence): the conifold instance
  is verified via the identification $\Theta_V \in \mathfrak{g}_\Gamma[V]$
  with the KS integration of $\CoHA(\bY)$; the general compact case
  is open.
\end{itemize}
\end{heal}

%======================================================================
\section{The grand atlas, rectified}
\label{sec:atlas}
%======================================================================

The atlas of \S1541--1591 of working\_notes is here rewritten with
(a) explicit subscripts on every $\kappa_\bullet$; (b) theorem/conjecture
status; (c) the Vol~III four-face spectrum where applicable;
(d) the direct-McKay source for the conifold. This is the
reader-facing atlas after Attacks 3 and 4.

\begin{center}
\small
\renewcommand{\arraystretch}{1.25}
\begin{longtable}{lccll}
\toprule
\textbf{CY3 family} & $\chi_\mathrm{top}$ & $\kch$ & \textbf{Status} & \textbf{Source / note} \\
\midrule
\endfirsthead
\toprule
\textbf{CY3 family} & $\chi_\mathrm{top}$ & $\kch$ & \textbf{Status} & \textbf{Source / note} \\
\midrule
\endhead
$\C^3$                         & ---    & $1$     & theorem        & $\kch(H_1) = 1$ (direct) \\
Conifold ($\cO(-1) \oplus \cO(-1) \to \bP^1$)
                               & ---    & $1$     & theorem        & direct McKay: $\CoHA(\bY)$ rank-1 per charge \\
Local $\bP^2$                  & ---    & $3/2$   & theorem        & McKay $\Z_3$ quiver; $\chi(\bP^2)/2$ \\
Local $\bP^1 \times \bP^1$     & ---    & $2$     & theorem        & toric web; $= h^{1,1}$ coincidence \\
Local $F_n$ ($n = 0,1,2$)      & ---    & $2$     & theorem        & Hirzebruch-surface McKay \\
Banana                         & $0$    & $0$     & conjecture     & BKY; class $\ge \mathbf{G}$ \\
$K3 \times E$                  & $0$    & $3$     & theorem        & Hodge supertrace; four-face $\{3,0,5,24\}$ \\
Enriques $\times E$            & $0$    & $4$     & theorem        & Allcock weight-4; $\kBKM = 4$ \\
Borcea--Voisin (K3 type)       & var    & $\chi/24$ & theorem      & involution CY; $\kBKM$ via base-form \\
Quintic $X_5 \subset \bP^4$    & $-200$ & $-25/3$ & theorem       & $\chi/24$; no BKM (non-integer) \\
$\bP^5[3,3]$                   & $-144$ & $-6$    & theorem        & $\chi/24$ \\
$\bP^5[2,4]$                   & $-176$ & $-22/3$ & theorem        & $\chi/24$ \\
$\bP^5[2,2,3]$                 & $-144$ & $-6$    & theorem        & $\chi/24$ \\
$\bP^7[2,2,2,2]$               & $-128$ & $-16/3$ & theorem        & $\chi/24$ \\
Schoen ($E \times_{\bP^1} E$)  & $0$    & $0$     & theorem        & Hodge supertrace \\
$G(2,5) \cap \{f_1, f_2\}$     & $-100$ & $-25/6$ & theorem        & Grassmannian CY \\
K3-fibered (general)           & var    & var     & conjecture     & Theorem~\ref{wn:thm:chi-neq-kappa} says $\kch \ne \chi/24$ for nontrivial K3-fibre\\
\bottomrule
\end{longtable}
\end{center}

\begin{observation}[Atlas patterns; rectified]
\label{obs:atlas-rectified}
\begin{enumerate}[label=(\roman*), leftmargin=1.2em]
\item Unsubscripted $\kappa_\bullet$ is forbidden; every atlas entry is $\kch$. The
  four-face spectrum $\{\kch, \kcat, \kBKM, \kfib\}$ is defined
  only for K3-modular families (K3 $\times$ E, Enriques $\times$ E,
  Borcea--Voisin with K3 quotient), where a Siegel-modular
  denominator form is known.
\item The BCOV $\chi/24$ identity holds for compact simply-connected
  CICYs with $h^{1,0} = 0$; it \emph{fails} for K3-fibered families
  (Theorem~\ref{wn:thm:chi-neq-kappa}); it is \emph{trivial} for
  families with $\chi = 0$ (K3 $\times$ E, Schoen, banana, Borcea
  in the K3 type).
\item The toric bound $r_{\max} \le N$ (for $N$ the vertex count of
  the toric web) is a \emph{conjecture}, not a theorem. Verified for
  $\C^3$ ($N = 1$, $r_{\max} = 1$), conifold ($N = 2$, $r_{\max} = 2$),
  local $\bP^2$ ($N = 3$, $r_{\max} = \infty$ in the M class). The
  bound fails in the M class; it applies only to G-class families.
\item All compact CY3 with $\chi_\mathrm{top} \ne 0$ are class
  $\mathbf{M}$ (Theorem~\ref{wn:thm:chi-neq-kappa} corollary).
\item A genuine BKM algebra requires the presence of a
  weight-$k$ paramodular form in the denominator formula; this
  picks out the K3-modular families above. Non-K3-modular
  families (quintic etc.) have no BKM realisation at present; the
  automorphic lift is open (open question 3 of working\_notes
  \S\ref{sec:updated-open}).
\end{enumerate}
\end{observation}

\paragraph{Status audit.}
Every ``theorem'' entry above refers to a proved $\kch$ value, not a
proved $A_X$ construction or a proved $\Phi$-identification. At
$d = 3$, $A_X$ is \emph{conjectural} (AP-CY6), and the category-level
$\Phi(D^b(X))$-output is conditional on CY-A$_3$ morphism-level
extension. The atlas records $\kch$ invariants only.

%======================================================================
\section{CY$_4$ prediction: audit and rectification}
\label{sec:cy4}
%======================================================================

Working\_notes \S\ref{sec:cy4} (line 2527) asserts the pattern
$\mathrm{CY}_d \to E_{d-2}$ predicts $\mathrm{CY}_4 \to \Etwo$, with
three supporting bullets.

\begin{attack}[CY$_4$ heuristic at $d = 4$]
\label{atk:6}
\textbf{Ghost theorem}: the $\Eone \hookrightarrow \Etwo$ enhancement
for CY$_4$ is visible at the level of $\CoHA$ for $K3 \times K3$
(factorisation $= $ two commuting K3-CoHAs, each $\Eone$, assembled
into $\Etwo$ by Dunn additivity).

\textbf{Precise error}: the Dunn additivity argument is stated as if
it closes the CY$_4 \to \Etwo$ prediction, but it assumes (i) the
K3 $\times$ K3 category decomposes as an $\Eone$-tensor product of
two K3 categories, (ii) the factorisation preserves the CY structure,
and (iii) $\CoHA(K3 \times K3) = \CoHA(K3) \otimes \CoHA(K3)$ as
$\Eone$-algebras with commuting structure. Assumption (iii) is
\emph{strong}: SV's CoHA of K3 surfaces (cohomological Hall of
Schiffmann--Vasserot) is not known to satisfy a clean external-tensor
decomposition; Maulik--Okounkov $R$-matrices for $K3 \times K3$ are
nontrivial and obstruct a trivial Dunn argument. The degree-4
potential $W$ in the bullet ``$d = 4$ potential $W$ has degree 4 vs 3''
is not the correct CY$_4$ potential for compact K3 $\times$ K3 (which
has $\cO_{K3 \times K3} \cong \wedge^4 T^*$ and no critical locus
structure); it is the local-CY$_4$ Jacobi potential, a different setup.

\textbf{Correct relation}: the CY$_4$ prediction is a \emph{genuine
heuristic} (Costello--Gaiotto, Kontsevich, Ginzburg) but is not a
theorem for compact K3 $\times$ K3. It holds:
\begin{itemize}[leftmargin=1.1em]
\item For the local CY$_4$ Jacobi-algebra setup (toric hypersurface of
  degree 4), where the $\Etwo$ structure on the CoHA comes from the
  $\Etwo$ factorisation algebra of the holomorphic Chern--Simons
  8-dimensional theory (Costello--Li 2016).
\item For products $\text{CY}_2 \times \text{CY}_2$ at the level of
  $\Eone$-bimodule data, with the $\Etwo$-enhancement conditional on
  commutativity of the two $\Eone$-structures (the Dunn+
  Eckmann--Hilton argument needs checking).
\end{itemize}
\end{attack}

\begin{heal}[CY$_4$ prediction rewritten as a labelled heuristic]
\label{heal:6}
Replace ``The pattern $\mathrm{CY}_d \to E_{d-2}$ predicts $\mathrm{CY}_4
\to \Etwo$. The swarm confirms:'' (line 2529) by:
\begin{quote}
\emph{Heuristic pattern.} $\mathrm{CY}_d$-geometry assigns an
$E_{d-2}$-algebra via the functor $\Phi$ (conjectural at $d \ge 3$,
proved at $d = 2$ as Vol~III Thm.~CY-A$_2$). The $d = 4$ case
predicts an $\Etwo$-algebra.

\emph{Partial confirmations.}
\begin{itemize}[leftmargin=1.1em]
\item \textit{Local CY$_4$ Jacobi algebras.} Holomorphic Chern--Simons
  on an 8-real-dimensional CY$_4$ gives an $\Etwo$-factorisation
  algebra at the level of Costello--Li.
\item \textit{$K3 \times K3$ via Dunn, \emph{conditional}.} If
  $\CoHA(K3 \times K3)$ admits a tensor decomposition
  $\CoHA(K3) \boxtimes \CoHA(K3)$ of $\Eone$-algebras with commuting
  structure, Dunn additivity produces an $\Etwo$-structure. The
  existence of the decomposition is \emph{open} (Maulik--Okounkov
  $R$-matrices on $K3 \times K3$ obstruct a trivial decomposition).
\item \textit{\v Cech descent.} The spectral sequence does not
  degenerate at $E_2$ for $\Etwo$-algebras;
  $\pi_1(\mathrm{Conf}_2(\R^2)) = \Z$ produces $E_2^{2, *} \ne 0$.
  This is consistent with the heuristic but not a verification.
\end{itemize}

The CY$_4 \to \Etwo$ pattern is \emph{conjectural} overall; the
compact case (K3 $\times$ K3, hyperkähler fourfolds) is open.
\end{quote}
\end{heal}

%======================================================================
\section{Required edits to working\_notes.tex}
\label{sec:edits}
%======================================================================

The edits below are stated as precise line-anchored interventions.

\begin{enumerate}[label=\textbf{E\arabic*.}, leftmargin=1.8em]

\item \textbf{Lines 414--416 (wall-crossing remark).} Replace with the
content of Heal~\ref{heal:1}: explicit dgLa, explicit $\Theta_\mathrm{I},
\Theta_\mathrm{II}$, explicit $\lambda$, explicit BCH height statement
with Reineke-convention caveat.

\item \textbf{Lines 1429--1452 (conifold subsection).} Rename
``(large-volume and flopped)'' on line 1432 to ``(Szendr\H{o}i NCDT and
large-volume DT)'' (Heal~\ref{heal:2}). Add the two-wall stratification
paragraph separating stability walls from the birational flop.

\item \textbf{Lines 1541--1591 (grand atlas).} Rewrite atlas body per
\S\ref{sec:atlas} and Heal~\ref{heal:3}, Heal~\ref{heal:4}. Conifold
source column becomes ``direct McKay''; $K3 \times E$ row expanded to
four-face spectrum $\{3, 0, 5, 24\}$; Enriques $\times E$ row gets
explicit $\kBKM = 4$.

\item \textbf{Lines 1802--1813 (updated open questions).} Items 4, 5,
13 updated per Heal~\ref{heal:5}.

\item \textbf{Lines 2529--2534 (CY$_4$ prediction).} Replace by the
labelled-heuristic statement of Heal~\ref{heal:6}; the three bullets
become ``partial confirmations'', with the compactness caveat and the
local-Jacobi clarification.

\end{enumerate}

%======================================================================
\section{Cache append: new anti-patterns}
\label{sec:cache}
%======================================================================

Append to \texttt{notes/antipatterns\_catalogue.md} (new entries).

\begin{itemize}[leftmargin=1.2em]

\item \textbf{AP-CY-WC1 -- dgLa unnamed (Critical).} \emph{Symptom}:
  stating ``wall-crossing is a gauge transformation of the MC element
  $\exp(\ad_\alpha)(\Theta_\mathrm{I}) = \Theta_\mathrm{II}$'' without
  specifying $(L, d, [,])$ or which $L^1$ contains $\Theta$.
  \emph{Counter}: every MC statement names the dgLa. For KS
  wall-crossing, the dgLa is $\mathfrak{g}_\Gamma[V]$ (zero differential;
  degree-zero concentrated; KS $(-1)^{\langle \cdot, \cdot \rangle}$
  signs); $\Theta = \sum_\gamma \Omega(\gamma) e_\gamma$;
  $\lambda = \log(\text{single-wall motivic product})$.

\item \textbf{AP-CY-WC2 -- Birational wall vs stability wall (High).}
  \emph{Symptom}: ``the two DT chambers (large-volume and flopped)'',
  conflating the Kähler (flop) wall with a $\Stab(D^b(X))$ stability
  wall. \emph{Counter}: flop is a derived equivalence $D^b(X) \to
  D^b(X^+)$ inducing an $\mathfrak{g}_\Gamma$ isomorphism; stability
  walls live inside a fixed $D^b(X)$. Use ``Szendr\H{o}i NCDT $
  \leftrightarrow$ large-volume DT'' for the conifold stability wall.

\item \textbf{AP-CY-WC3 -- Pentagon at height 3 (Medium).} \emph{Symptom}:
  asserting the BCH pentagon holds at all orders in the Lie-theoretic
  semi-classical dgLa. \emph{Counter}: it holds through height 2 in the
  Lie-theoretic dgLa with KS signs, and to all orders only in the
  quantum-torus dgLa or at the motivic Hall-algebra level (Gross--
  Siebert scattering-diagram consistency; \texttt{arXiv:0911.2780}).
  Working\_notes \S2547 is the canonical statement; any Lie-algebra
  verification of the pentagon must name the height.

\item \textbf{AP-CY-WC4 -- ``DT genus-$n$'' ambiguity (Medium).}
  \emph{Symptom}: labelling an invariant ``DT genus-$n$'' without
  specifying whether this refers to (i) the Gopakumar--Vafa genus-$n$
  free energy $F_n$, (ii) a genus-$n$ bar-complex $B_\Eone^n$, or
  (iii) modular-weight $n$ of the DT partition function.
  \emph{Counter}: use ``direct McKay'', ``$F_1^{\mathrm{GV}}$'',
  ``bar $B^1$'', or ``modular weight 1'' as appropriate.
  $\kch(\text{conifold}) = 1$ is from direct McKay (per CLAUDE.md),
  not from $F_1^{\mathrm{GV}}$.

\item \textbf{AP-CY-WC5 -- $\{\kch, \kcat, \kBKM, \kfib\}$ elision in
  the atlas (High).} \emph{Symptom}: atlas rows for K3-modular CY3
  families quoting only $\kch$ or only $\kBKM$, inviting the reader
  to infer $\kBKM = \kch + \kfib$ or similar false identity.
  \emph{Counter}: for K3-modular families ($K3 \times E$, Enriques
  $\times$ E, Borcea--Voisin K3 type), display the full four-face
  spectrum; for non-K3-modular families, display only $\kch$ and
  note that no BKM realisation is known.

\end{itemize}

%======================================================================
\section{Summary}
%======================================================================

The five adversarial attacks identified five distinct gaps in the
wall-crossing / MC gauge equivalence / conifold / atlas / CY$_4$
apparatus:
(1) unnamed dgLa at \S\ref{wn:sec:wall-crossing-mc};
(2) Kähler/stability-wall conflation at \S\ref{subsec:conifold-full};
(3) ambiguous ``DT genus-1'' source column for $\kch(\text{conifold})$;
(4) atlas $K3 \times E$ row missing three of four $\kappa_\bullet$ entries;
(5) CY$_4 \to \Etwo$ presented as near-theorem when it is heuristic.

The corresponding heals supplied: (1) the explicit dgLa
$\mathfrak{g}_\Gamma[V]$ with $\Theta_\mathrm{I} = e_{\gamma_1} +
e_{\gamma_2}$, $\Theta_\mathrm{II} = e_{\gamma_1} + e_{\gamma_1 + \gamma_2}
+ e_{\gamma_2}$, $\lambda = \tfrac{1}{2}e_{\gamma_1 + \gamma_2}$;
(2) the two-wall stratification with flop as Bridgeland equivalence;
(3) ``direct McKay'' replacing ``DT genus-1'' per the CLAUDE.md cache;
(4) the four-face spectrum $\{3, 0, 5, 24\}$ for $K3 \times E$;
(5) CY$_4 \to \Etwo$ relabelled as heuristic with local-Jacobi and
$K3 \times K3$-conditional partial confirmations.

Five new AP-CY-WC entries appended to the anti-pattern catalogue;
reader-facing line edits to working\_notes.tex listed in
\S\ref{sec:edits}.

\end{document}
